% =====================================================================
%  Informe Etapa 1 -- Trabajo práctico de Econometría II
%  Esqueleto con el diseño de apuntes-latex (standalone, sin cajas).
%  Las partes en gris entre corchetes se reemplazan por el contenido.
% =====================================================================
\documentclass[11pt,letterpaper]{article}

% --- Idioma y fuentes ------------------------------------------------
\usepackage[spanish,es-noshorthands,es-tabla]{babel}
\usepackage[T1]{fontenc}
\usepackage{lmodern}

% --- Geometría y encabezado ------------------------------------------
\usepackage[a4paper,margin=2.5cm]{geometry}
\usepackage{fancyhdr}
\setlength{\headheight}{38.22325pt}
\addtolength{\topmargin}{-6.22325pt}

% --- Matemáticas ------------------------------------------------------
\usepackage{amsmath,amssymb}

% --- Colores, tablas, listas, código ----------------------------------
\usepackage[dvipsnames,table]{xcolor}
\usepackage{booktabs}
\usepackage{array}
\usepackage{longtable}
\usepackage{enumitem}
\usepackage{graphicx}
\usepackage{listings}
\renewcommand{\arraystretch}{1.5}

\lstset{
  language=R,
  basicstyle=\ttfamily\footnotesize,
  commentstyle=\color{gray}\itshape,
  keywordstyle=\color{blue!70!black},
  stringstyle=\color{teal!60!black},
  breaklines=true,
  columns=fullflexible,
  showstringspaces=false
}

% Texto pendiente de completar (se borra al rellenar)
\newcommand{\completar}[1]{\textcolor{gray}{\textit{[#1]}}}

\pagestyle{fancy}
\fancyhf{}
\fancyhead[L]{\textbf{Econometría II}\\ Trabajo práctico --- Etapa 1}
\fancyhead[C]{Universidad de Santiago de Chile}
% Logo: sube logo.png.png a Overleaf y descomenta la línea de abajo
% \fancyhead[R]{\includegraphics[height=1.2cm]{logo.png.png}}
\fancyfoot[C]{\thepage}

\begin{document}

% =====================================================================
%  PORTADA
% =====================================================================
\begin{center}
  {\LARGE \textbf{Etapa 1: construcción de la base de datos}}\\[4pt]
  {\large Econometría II --- Admisión 2026}
\end{center}
\vspace{0.5em}

\begin{center}
\begin{tabular}{ll}
\toprule
\textbf{Integrantes} & \completar{Nombre 1} \\
                     & \completar{Nombre 2} \\
                     & \completar{Nombre 3} \\
\midrule
\textbf{Profesor}    & \completar{Nombre del profesor} \\
\textbf{Fecha}       & \completar{Fecha de entrega} \\
\bottomrule
\end{tabular}
\end{center}

\bigskip

% =====================================================================
\subsection*{1 \; Objetivo}
% =====================================================================

\completar{Qué se busca con la base: un solo data frame con los
matriculados del proceso 2026, su PAES, sus notas de enseñanza media y
variables del colegio de procedencia, para estimar en la Etapa 3 la
probabilidad de matricularse dentro de las tres primeras preferencias.}

\medskip\hrule\medskip

% =====================================================================
\subsection*{2 \; Fuentes de datos}
% =====================================================================

\completar{Origen: portal de bases de datos del DEMRE, proceso de admisión
2026. Llave de unión entre archivos: \texttt{ID\_aux}.}

\begin{center}
\begin{tabular}{>{\raggedright\arraybackslash}p{4.2cm}
                >{\raggedright\arraybackslash}p{5.8cm}
                >{\raggedleft\arraybackslash}p{2.6cm}}
\toprule
\textbf{Archivo} & \textbf{Contenido} & \textbf{Filas} \\
\midrule
\completar{Archivo B} & \completar{Inscritos} & \completar{n} \\
\completar{Archivo C} & \completar{Rinden PAES} & \completar{n} \\
\completar{Archivo D} & \completar{Postulaciones} & \completar{n} \\
\completar{Archivo de matrícula} & \completar{Matriculados} & \completar{n} \\
\bottomrule
\end{tabular}
\end{center}

\medskip\hrule\medskip

% =====================================================================
\subsection*{3 \; Construcción de la base}
% =====================================================================

\subsubsection*{3.1 \; Selección de columnas y unión de archivos}

\completar{Qué columnas se tomaron de cada archivo y cómo se unieron
(tipo de \emph{join} y archivo de partida).}

\subsubsection*{3.2 \; Filtros aplicados}

\completar{Un filtro por ítem, con la variable y el valor que se usó.}

\begin{enumerate}[leftmargin=1.5em, itemsep=0.3em, label=\textbf{\alph*)}]
  \item \textbf{Vía y tipo de matrícula:} \completar{matriculados por vía
        regular con tipo de matrícula «postulación regular».}
  \item \textbf{Preferencia:} \completar{si se matriculó, que haya sido en
        una preferencia hasta el tercer lugar.}
  \item \textbf{Promoción del año:} \completar{rindieron la PAES siendo
        alumnos regulares de la promoción 2025.}
\end{enumerate}

\subsubsection*{3.3 \; Observaciones en cada paso}

\begin{center}
\begin{tabular}{>{\raggedright\arraybackslash}p{9cm}
                >{\raggedleft\arraybackslash}p{3cm}}
\toprule
\textbf{Paso} & \textbf{Observaciones} \\
\midrule
\completar{Base inicial} & \completar{n} \\
\completar{Filtro a)} & \completar{n} \\
\completar{Filtro b)} & \completar{n} \\
\completar{Filtro c)} & \completar{n} \\
\rowcolor{yellow!30}
\textbf{Base final} & \completar{n} \\
\bottomrule
\end{tabular}
\end{center}

\medskip\hrule\medskip

% =====================================================================
\subsection*{4 \; Variable de respuesta}
% =====================================================================

\completar{Definición elegida por el grupo y por qué.}

\[
Y_i =
\begin{cases}
1 & \text{si } \completar{condición} \\[4pt]
0 & \text{si } \completar{condición}
\end{cases}
\]

\begin{center}
\begin{tabular}{lrr}
\toprule
\textbf{Valor de $Y$} & \textbf{Estudiantes} & \textbf{Porcentaje} \\
\midrule
$Y = 1$ & \completar{n} & \completar{\%} \\
$Y = 0$ & \completar{n} & \completar{\%} \\
\midrule
\textbf{Total} & \completar{n} & 100\% \\
\bottomrule
\end{tabular}
\end{center}

\medskip\hrule\medskip

% =====================================================================
\subsection*{5 \; Variables creadas}
% =====================================================================

\subsubsection*{5.1 \; Variables del establecimiento}

\completar{Tamaño, matrícula, etc. Cómo se agregaron por colegio
(\texttt{RBD}) y con qué grupo de estudiantes se calcularon, para que sea
consistente con un corte transversal del proceso 2026.}

\subsubsection*{5.2 \; Variables calculadas a partir de las existentes}

\completar{Por ejemplo: colegio rural o no, alumnos por curso, dummies.
Fórmula o regla de cada una.}

\medskip\hrule\medskip

% =====================================================================
\subsection*{6 \; Datos sin información}
% =====================================================================

\completar{Valores faltantes y códigos especiales (por ejemplo, 99 =
sin información) y qué se hizo con cada uno.}

\begin{center}
\begin{tabular}{>{\raggedright\arraybackslash}p{3.8cm}
                >{\raggedleft\arraybackslash}p{2.4cm}
                >{\raggedright\arraybackslash}p{6.4cm}}
\toprule
\textbf{Variable} & \textbf{Casos} & \textbf{Tratamiento} \\
\midrule
\completar{variable} & \completar{n} & \completar{qué se hizo} \\
\completar{variable} & \completar{n} & \completar{qué se hizo} \\
\bottomrule
\end{tabular}
\end{center}

\medskip\hrule\medskip

% =====================================================================
\subsection*{7 \; Definición de las variables}
% =====================================================================

\completar{Diccionario de la base final: una fila por variable.}

\begin{longtable}{>{\raggedright\arraybackslash}p{3.6cm}
                  >{\raggedright\arraybackslash}p{1.6cm}
                  >{\raggedright\arraybackslash}p{1.8cm}
                  >{\raggedright\arraybackslash}p{5.4cm}}
\toprule
\textbf{Variable} & \textbf{Fuente} & \textbf{Tipo} & \textbf{Definición} \\
\midrule
\endfirsthead
\toprule
\textbf{Variable} & \textbf{Fuente} & \textbf{Tipo} & \textbf{Definición} \\
\midrule
\endhead
\bottomrule
\endfoot
\multicolumn{4}{l}{\textit{Identificación y respuesta}} \\
\texttt{ID\_aux} & \completar{archivo} & \completar{tipo} & \completar{definición} \\
\texttt{Y} & \completar{creada} & \completar{dummy} & \completar{definición} \\
\midrule
\multicolumn{4}{l}{\textit{Estudiante}} \\
\completar{variable} & \completar{archivo} & \completar{tipo} & \completar{definición} \\
\completar{variable} & \completar{archivo} & \completar{tipo} & \completar{definición} \\
\midrule
\multicolumn{4}{l}{\textit{PAES y notas}} \\
\completar{variable} & \completar{archivo} & \completar{tipo} & \completar{definición} \\
\completar{variable} & \completar{archivo} & \completar{tipo} & \completar{definición} \\
\midrule
\multicolumn{4}{l}{\textit{Establecimiento}} \\
\completar{variable} & \completar{archivo} & \completar{tipo} & \completar{definición} \\
\completar{variable} & \completar{archivo} & \completar{tipo} & \completar{definición} \\
\end{longtable}

\medskip\hrule\medskip

% =====================================================================
\subsection*{8 \; Base final}
% =====================================================================

\completar{Dimensiones (filas y columnas), una fila por estudiante, sin
datos faltantes, nombre del archivo guardado.}

% =====================================================================
\newpage
\subsection*{Anexo \; Código en R}
% =====================================================================

\completar{Pegar el código dentro del bloque de abajo. Otra opción en
Overleaf: subir el script y reemplazar el bloque por}

\begin{center}
\texttt{\textbackslash lstinputlisting\{04-base-grupo.R\}}
\end{center}

\begin{lstlisting}
# codigo del grupo
\end{lstlisting}

\end{document}
